%%%%%%%%%%%%%%%%%%%%%%%%%%%%%%%%%%%%%%%%%%%%%%%%%%%%%%%%%%%%%%%%%%%%%%%%%%%%%%%%

\part{Elementary quantum algorithms}
\begin{frame}
  \partpage
  \label{part:elementary}
\end{frame}

%%%%%%%%%%%%%%%%%%%%%%%%%%%%%%%%%%%%%%%%%%%%%%%%%%%%%%%%%%%%%%%%%%%%%%%%%%%%%%%%

\begin{frame}{Black box problems}

Standard computational problem: determine a property of some input data
\begin{itemize}
  \item Example: Find the prime factors of $N$
\end{itemize}

\pause
\bigskip

Alternate model: Input is provided by a \emph{black box} (or \emph{oracle})
\begin{itemize}
  \item Query: On input $x$, black box returns $f(x)$ \pause
  \item Determine a property of $f$ using as few queries as possible
  \item The minimum number of queries is the \emph{query complexity} \pause
  \item Example: Given a black box for $f:\{1,2,\ldots,N\} \to \{0,1\}$,
    is there some $x$ such that $f(x)=1$? \pause
  \item Why black boxes? \pause
  \begin{itemize}
    \item Facilitates proving lower bounds \pause
    \item Can lead to algorithms for standard problems
  \end{itemize}
\end{itemize}

\end{frame}

%%%%%%%%%%%%%%%%%%%%%%%%%%%%%%%%%%%%%%%%%%%%%%%%%%%%%%%%%%%%%%%%%%%%%%%%%%%%%%%%

\begin{frame}{Black boxes for reversible/quantum computing}

Black box ~~
    \raisebox{.35em}{\Qcircuit @C=1em @R=.7em {
  & \lstick{x} & \gate{f} & \rstick{f(x)} \qw \\
  }} \hspace{2em}
~~ is not reversible
  
  \pause
  \bigskip
  
Reversible version: ~~
    \raisebox{.35em}{\Qcircuit @C=1em @R=.7em {
  & \lstick{x} & \multigate{1}{f} & \rstick{x} \qw \\
  & \lstick{z} & \ghost{f} & \rstick{z \oplus f(x)} \qw \\
  }}
  
  \pause
  \bigskip
  
Given a circuit that computes $f$ non-reversibly, we can implement the reversible version with little overhead

  \pause
  \bigskip
  
Quantum version: ~~
    \raisebox{.35em}{\Qcircuit @C=1em @R=.7em {
  & \lstick{|x\>} & \multigate{1}{f} & \rstick{|x\>} \qw \\
  & \lstick{|z\>} & \ghost{f} & \rstick{|z \oplus f(x)\>} \qw \\
  }}

  \pause
  \bigskip

A reversible circuit is a quantum circuit

\end{frame}

%%%%%%%%%%%%%%%%%%%%%%%%%%%%%%%%%%%%%%%%%%%%%%%%%%%%%%%%%%%%%%%%%%%%%%%%%%%%%%%%

\begin{frame}{Deutsch's problem}

\begin{prob}
\begin{itemize}
\item Given: a black-box function $f:\{0,1\} \to \{0,1\}$
\item Task: determine whether $f$ is \emph{constant} or \emph{balanced}
\end{itemize}
\end{prob}

\pause

\begin{center}
\begin{tabular}{cccc}
  \begin{tabular}{c|c}$x$ & $f_1(x)$ \\ \hline 0 & 0 \\ 1 & 0\end{tabular} &
  \begin{tabular}{c|c}$x$ & $f_2(x)$ \\ \hline 0 & 1 \\ 1 & 1\end{tabular} &
  \begin{tabular}{c|c}$x$ & $f_3(x)$ \\ \hline 0 & 0 \\ 1 & 1\end{tabular} &
  \begin{tabular}{c|c}$x$ & $f_4(x)$ \\ \hline 0 & 1 \\ 1 & 0\end{tabular} \\[-4pt]
  \multicolumn{2}{@{}c@{}}{\onslide<3->{$\underbrace{\hspace{1.6in}}_{\mbox{constant: $f(0)=f(1)$}}$}} &
  \multicolumn{2}{@{}c@{}}{\onslide<4->{$\underbrace{\hspace{1.6in}}_{\mbox{balanced: $f(0)\ne f(1)$}}$}}
\end{tabular}
\end{center}

\bigskip

\onslide<5->{How many queries are needed?}
\begin{itemize}
\item<6-> Classically: 2 queries are necessary and sufficient \pause
\item<7-> Quantumly: ?
\end{itemize}

\end{frame}

%%%%%%%%%%%%%%%%%%%%%%%%%%%%%%%%%%%%%%%%%%%%%%%%%%%%%%%%%%%%%%%%%%%%%%%%%%%%%%%%

\begin{frame}{Toward a quantum algorithm for Deutsch's problem}

Quantum black box for $f$: \hspace{1ex}
  \raisebox{.35em}{\Qcircuit @C=1em @R=.7em {
  & \lstick{|x\>} & \multigate{1}{f} & \rstick{|x\>} \qw \\
  & \lstick{|z\>} & \ghost{f} & \rstick{|z \oplus f(x)\>} \qw \\
  }}

\bigskip\bigskip

\onslide<2->{Compute $f$ in superposition: \hspace{1ex}
  \raisebox{.35em}{
  \Qcircuit @C=1em @R=.7em {
  & \lstick{|0\>} & \gate{\color<3|handout:0>{red} H} & \multigate{1}{\color<4|handout:0>{red} f} & \qw \\
  & \lstick{|0\>} & \qw & \ghost{f} & \qw \\
  }}
\medskip
}

\onslide<3->{
\begin{align*}
  |0\>\ot|0\>
  &{\color<3|handout:0>{red} \mapsto \frac{|0\>+|1\>}{\sqrt2}\ot|0\>} \\
\onslide<4->{
  &{\color<4|handout:0>{red} \mapsto \frac{1}{\sqrt2}(|0\>\ot|f(0)\>+|1\>\ot|f(1)\>)}
  }
\end{align*}
}

\onslide<5->{
Can't extract more than one bit of information about $f$
}

\end{frame}


%%%%%%%%%%%%%%%%%%%%%%%%%%%%%%%%%%%%%%%%%%%%%%%%%%%%%%%%%%%%%%%%%%%%%%%%%%%%%%%%

\begin{frame}{Phase kickback}

% Recall phase kickback example from lecture 1:
%   \[
%   \Qcircuit @C=1em @R=1.5em {
%   & \lstick{|x\>} & \ctrl{1} & \rstick{e^{2\pi\ii\varphi x}|x\>} \qw \\
%   & \lstick{|\psi\>} & \gate{U} & \rstick{|\psi\>} \qw \\
%   }
%   \]
% 
%   \bigskip
% 
% \onslide<2->{Similar trick for a black box:

\bigskip

Quantum black box for $f$: \hspace{1ex}
  \raisebox{.35em}{\Qcircuit @C=1em @R=.7em {
  & \lstick{|x\>} & \multigate{1}{f} & \rstick{|x\>} \qw \\
  & \lstick{|z\>} & \ghost{f} & \rstick{|z \oplus f(x)\>} \qw \\
  }}
  
\bigskip\bigskip

\onslide<2->{
Phase kickback:
  \[
  \only<2-6|handout:0>{
  \Qcircuit @C=1em @R=.7em {
  & \lstick{\color<3>{red}{|x\>}} & \multigate{1}{\color<5>{red}{f}} & \qw \\
  & \lstick{\color<3>{red}{\frac{|0\>-|1\>}{\sqrt2}}} & \ghost{f} & \qw \\
  }}
  \only<7->{
  \Qcircuit @C=1em @R=.7em {
  & \lstick{|x\>} & \multigate{1}{f} & \rstick{\color<7|handout:0>{red}{(-1)^{f(x)}|x\>}} \qw \\
  & \lstick{\frac{|0\>-|1\>}{\sqrt2}} & \ghost{f} & \rstick{\color<7|handout:0>{red}{\frac{|0\>-|1\>}{\sqrt2}}} \qw \\
  }}
  \]
}

\begin{align*}
\onslide<3->{
  \color<3|handout:0>{red}{|x\> \ot \smfrac{1}{\sqrt2} (|0\>-|1\>)}
  \onslide<4->{&=\smfrac{1}{\sqrt2}(|x\>\ot|0\> - |x\>\ot|1\>) \\}
  \onslide<5->{&\color<5|handout:0>{red}{\mapsto\smfrac{1}{\sqrt2}(|x\>\ot|f(x)\> - |x\>\ot|1\oplus f(x)\>)} \\}
  \onslide<6->{&=|x\>\ot \smfrac{1}{\sqrt2} (|f(x)\> - |\overline{f(x)}\>) \\}
 \onslide<7->{&=\color<7|handout:0>{red}{\only<7|handout:0>{(-1)^{f(x)}}\only<8>{\underbrace{(-1)^{f(x)}}_{\mathclap{\text{not necessarily global}}}}|x\>\ot\smfrac{1}{\sqrt2}(|0\> - |1\>)}
 % following line is just here for spacing
 \onslide<5|handout:0>{\underbrace{(-1)^{f(x)}}_{\mathclap{\text{not necessarily global}}}}
 \\}
}
\end{align*}

\end{frame}

%%%%%%%%%%%%%%%%%%%%%%%%%%%%%%%%%%%%%%%%%%%%%%%%%%%%%%%%%%%%%%%%%%%%%%%%%%%%%%%%

\begin{frame}{Quantum algorithm for Deutsch's problem}

\[
  \Qcircuit @C=1em @R=.7em {
  & \lstick{|0\>} & \gate{\color<2|handout:0>{red}H} & \multigate{1}{\color<3|handout:0>{red}f} & \gate{\color<5|handout:0>{red}H} & \meter & \rstick{\onslide<6->{\color<6|handout:0>{red}{f(0) \oplus f(1)}}} \cw \\
  & \lstick{\frac{|0\>-|1\>}{\sqrt2}} & \qw & \ghost{f} & \qw & \qw \\
  }
\]

\bigskip

\onslide<2->{
\begin{align*}
  |0\>\otimes\frac{|0\>-|1\>}{\sqrt2}
  &\color<2|handout:0>{red}{\mapsto \frac{|0\>+|1\>}{\sqrt2} \otimes \frac{|0\>-|1\>}{\sqrt2}} \\
\onslide<3->{
  &\color<3|handout:0>{red}{\mapsto \frac{(-1)^{f(0)}|0\>+(-1)^{f(1)}|1\>}{\sqrt2} \otimes \frac{|0\>-|1\>}{\sqrt2}} \\
}
\onslide<4->{
  &= (-1)^{f(0)}\frac{|0\>+(-1)^{f(0) \oplus f(1)}|1\>}{\sqrt2} \otimes \frac{|0\>-|1\>}{\sqrt2} \\
}
\onslide<5->{
  &\color<5|handout:0>{red}{\mapsto (-1)^{f(0)} |{\color<6|handout:0>{red} f(0) \oplus f(1)}\> \otimes \frac{|0\>-|1\>}{\sqrt2}}
}
\end{align*}
}

\vspace{-2ex}

\onslide<7->{
\begin{center}
1 quantum query vs.\ 2 classical queries!
\end{center}
}

\end{frame}

%%%%%%%%%%%%%%%%%%%%%%%%%%%%%%%%%%%%%%%%%%%%%%%%%%%%%%%%%%%%%%%%%%%%%%%%%%%%%%%%

\begin{frame}{The Deutsch-Jozsa problem}

\begin{prob}
\begin{itemize}
\item Given: a black-box function $f:\{0,1\}^n \to \{0,1\}$
\item Promise: $f$ is either
\begin{tabular}{@{}ll@{}}
  \emph{constant} & ($f(x)$ is independent of $x$) \\
  or \emph{balanced} & ($f(x)=0$ for exactly half the values of $x$)
\end{tabular}
\item Task: determine whether $f$ is {constant} or {balanced}
\end{itemize}
\end{prob}

\pause
\bigskip

How many queries are needed? \pause
\begin{itemize}
\item Classically: $2^n/2 + 1$ queries to answer with certainty \pause
\item Quantumly: ?
\end{itemize}

\end{frame}

%%%%%%%%%%%%%%%%%%%%%%%%%%%%%%%%%%%%%%%%%%%%%%%%%%%%%%%%%%%%%%%%%%%%%%%%%%%%%%%%

\begin{frame}{Phase kickback for a Boolean function of $n$ bits}

Black box function: \hspace{1.5ex}
  \raisebox{.35em}{\Qcircuit @C=1em @R=.7em {
  & \lstick{|x_1\>} & \multigate{3}{f} & \rstick{|x_1\>} \qw \\
  & \raisebox{4pt}{\vdots} & \push{\rule{0em}{1em}} & \raisebox{4pt}{\vdots} \\
  & \lstick{|x_n\>} & \ghost{f} & \rstick{|x_n\>} \qw \\
  & \lstick{|z\>} & \ghost{f} & \rstick{|z \oplus f(x)\>} \qw \\
  }}

  \bigskip
  \bigskip

\onslide<2->{Phase kickback: 
\[
  |x_1\>\ot\cdots\ot|x_n\>\ot\frac{|0\>-|1\>}{\sqrt2}
  \mapsto
  (-1)^{f(x)}
  |x_1\>\ot\cdots\ot|x_n\>\ot\frac{|0\>-|1\>}{\sqrt2}
\]
}

\end{frame}

%%%%%%%%%%%%%%%%%%%%%%%%%%%%%%%%%%%%%%%%%%%%%%%%%%%%%%%%%%%%%%%%%%%%%%%%%%%%%%%%

\begin{frame}{Quantum algorithm for the Deutsch-Jozsa problem}
  
  \[
  \Qcircuit @C=1em @R=.7em {
  & \lstick{|0\>} & \gate{\color<2|handout:0>{red}H} & \multigate{3}{\color<4|handout:0>{red}f} & \gate{\color<5|handout:0>{red}H} & \meter \\
  & & \raisebox{4pt}{\vdots} & \push{\rule{0em}{1em}} & \raisebox{4pt}{\vdots} & \raisebox{4pt}{\vdots} \\
  & \lstick{|0\>} & \gate{\color<2|handout:0>{red}H} & \ghost{f} & \gate{\color<5|handout:0>{red}H} & \meter \\
  & \lstick{\frac{|0\>-|1\>}{\sqrt2}} & \qw & \ghost{f} & \qw & \qw \\
  }
\]

\bigskip
\bigskip

\onslide<2->{
\begin{align*}
  |0\>^{\otimes n}\otimes\frac{|0\>-|1\>}{\sqrt2}
  &\color<2|handout:0>{red}{\mapsto \left(\frac{|0\>+|1\>}{\sqrt2}\right)^{\otimes n}  \otimes \frac{|0\>-|1\>}{\sqrt2}} \\
\onslide<3->{
  &= \frac{1}{\sqrt{2^n}} \sum_{x \in \{0,1\}^n} \! |x\> \otimes \frac{|0\>-|1\>}{\sqrt2} \\
}
\onslide<4->{
  &\color<4|handout:0>{red}{\mapsto \frac{1}{\sqrt{2^n}} \sum_{x \in \{0,1\}^n} \! (-1)^{f(x)} |x\> \otimes \frac{|0\>-|1\>}{\sqrt2}} \\
}
\end{align*}
}

\end{frame}

%%%%%%%%%%%%%%%%%%%%%%%%%%%%%%%%%%%%%%%%%%%%%%%%%%%%%%%%%%%%%%%%%%%%%%%%%%%%%%%%

\begin{frame}{Hadamard transform}

What do the final Hadamard gates do?

\onslide<2->{
\begin{align*}
  H|x\>
  &\onslide<3->{= \frac{1}{\sqrt2}(|0\> + (-1)^x|1\>)} \\
  &\onslide<4->{= \frac{1}{\sqrt2} \sum_{y \in \{0,1\}} (-1)^{xy} |y\>}
\end{align*}
}

\vspace{-2ex}

\onslide<5->{
\begin{align*}
  H^{\otimes n}(|x_1\> \otimes \cdots \otimes |x_n\>)
  &\onslide<6->{= \bigotimes_{i=1}^n \left(\frac{1}{\sqrt2} \sum_{y_i \in \{0,1\}} (-1)^{x_i y_i} |y_i\>\right)} \\
  &\onslide<7->{= \frac{1}{\sqrt{2^n}} \sum_{y \in \{0,1\}^n} (-1)^{x \cdot y} |y\>}
\end{align*}
}

\end{frame}

%%%%%%%%%%%%%%%%%%%%%%%%%%%%%%%%%%%%%%%%%%%%%%%%%%%%%%%%%%%%%%%%%%%%%%%%%%%%%%%%

\begin{frame}{Quantum D-J algorithm: Finishing up}
  
\begin{align*}
  \frac{1}{\sqrt{2^n}} \sum_{x \in \{0,1\}^n} \! (-1)^{f(x)} |x\>
  \onslide<2->{
  \stackrel{H^{\otimes n}}{\mapsto}
  \frac{1}{2^n} \sum_{x,y \in \{0,1\}^n} \! (-1)^{f(x)} (-1)^{x \cdot y} |y\>}
\end{align*}

\begin{itemize}
  \item<3-> If $f$ is constant, the amplitude of $|y\>$ is
  \[
    \pm \frac{1}{2^n}\sum_{x \in \{0,1\}^n} (-1)^{x \cdot y}
    \onslide<4->{= \pm \begin{cases}
    1 & \text{if $y=0\ldots0$} \\
    0 & \text{otherwise}
    \end{cases}}
  \]
  \onslide<4->{
  so we definitely measure $0\ldots0$
  }
  \item<5-> If $f$ is balanced, the amplitude of $|0\ldots0\>$ is
  \[
    \sum_{x \in \{0,1\}^n} (-1)^{f(x)} \onslide<6->{= 0}
  \]
  \onslide<6->{
  so we measure some nonzero string
  }
\end{itemize}

\end{frame}

%%%%%%%%%%%%%%%%%%%%%%%%%%%%%%%%%%%%%%%%%%%%%%%%%%%%%%%%%%%%%%%%%%%%%%%%%%%%%%%%

\begin{frame}{The Deutsch-Jozsa problem: Quantum vs.\ classical}

Above quantum algorithm uses only one query.

\bigskip

Need $2^n/2 + 1$ classical queries to answer with certainty.

\pause
\bigskip

What about randomized algorithms? \pause Success probability arbitrarily close to $1$ with a constant number of queries.

\pause
\bigskip
 
Can we get a separation between randomized and quantum computation?

\end{frame}

%%%%%%%%%%%%%%%%%%%%%%%%%%%%%%%%%%%%%%%%%%%%%%%%%%%%%%%%%%%%%%%%%%%%%%%%%%%%%%%%

\begin{frame}{Simon's problem}
  
\begin{prob}
\begin{itemize}
\item Given: a black-box function $f:\{0,1\}^n \to \{0,1\}^m$
\item Promise: there is some $s \in \{0,1\}^n$ such that $f(x)=f(y)$ if and only if $x=y$ or $x=y \oplus s$
\item Task: determine $s$
\end{itemize}
\end{prob}

\pause
\bigskip

One classical strategy:
\begin{itemize}
  \item query a random $x$
  \item repeat until we find $x_i \ne x_j$ such that $f(x_i)=f(x_j)$
  \item output $x_i \oplus x_j$
\end{itemize}

\pause
\bigskip

By the birthday problem, this uses about $\sqrt{2^n}$ queries.

\pause
\bigskip

It can be shown that this strategy is essentially optimal.

\end{frame}

%%%%%%%%%%%%%%%%%%%%%%%%%%%%%%%%%%%%%%%%%%%%%%%%%%%%%%%%%%%%%%%%%%%%%%%%%%%%%%%%

\begin{frame}{Quantum algorithm for Simon's problem}
  
Quantum black box: $|x\>\ot|y\> \mapsto |x\>\ot|y \oplus f(x)\>$ \smallskip\\
\hfill($x \in \{0,1\}^n$, $y \in \{0,1\}^m$)

\pause
\bigskip
  
\[
  \Qcircuit @C=1em @R=.7em {
  & \lstick{|0\>} & \gate{H} & \multigate{5}{f} & \gate{H} & \meter \\
  & & \raisebox{4pt}{\vdots} & \push{\rule{0em}{1em}} & \raisebox{4pt}{\vdots} & \raisebox{4pt}{\vdots} \\
  & \lstick{|0\>} & \gate{H} & \ghost{f} & \gate{H} & \meter \\
  & \lstick{|0\>} & \qw & \ghost{f} & \qw & \qw \\
  & \raisebox{4pt}{\vdots} & & \push{\rule{0em}{1em}} & & \raisebox{4pt}{\vdots} \\
  & \lstick{|0\>} & \qw & \ghost{f} & \qw & \qw \\
  }
\]

\bigskip

Repeat many times and post-process the measurement outcomes

\end{frame}

%%%%%%%%%%%%%%%%%%%%%%%%%%%%%%%%%%%%%%%%%%%%%%%%%%%%%%%%%%%%%%%%%%%%%%%%%%%%%%%%

\begin{frame}{Quantum algorithm for Simon's problem: Analysis I}

\begin{columns}
\begin{column}{5cm}
\[
  \Qcircuit @C=1em @R=.7em {
  & \lstick{|0\>} & \gate{\color<3|handout:0>{red}H} & \multigate{5}{\color<4|handout:0>{red}f} & \gate{H} & \meter \\
  & & \raisebox{4pt}{\vdots} & \push{\rule{0em}{1em}} & \raisebox{4pt}{\vdots} & \raisebox{4pt}{\vdots} \\
  & \lstick{|0\>} & \gate{\color<3|handout:0>{red}H} & \ghost{f} & \gate{H} & \meter \\
  & \lstick{|0\>} & \qw & \ghost{f} & \qw & \qw \\
  & \raisebox{4pt}{\vdots} & & \push{\rule{0em}{1em}} & & \raisebox{4pt}{\vdots} \\
  & \lstick{|0\>} & \qw & \ghost{f} & \qw & \qw \\
  }
\]
\end{column}
\begin{column}{5cm}
  \begin{align*}
  &\onslide<2->{|0\>^{\otimes n} \otimes |0\>^{\otimes m}} \\
  &\onslide<3->{\color<3|handout:0>{red}\mapsto \frac{1}{\sqrt{2^n}} \sum_{x \in \{0,1\}^n} |x\> \otimes |0\>^{\otimes m}} \\
  &\onslide<4->{\color<4|handout:0>{red}\mapsto \frac{1}{\sqrt{2^n}} \sum_{x \in \{0,1\}^n} |x\> \otimes |f(x)\>} \\
  &\onslide<5->{= \frac{1}{\sqrt{2^{n-1}}} \sum_{x \in R} \frac{|x\>+|x \oplus s\>}{\sqrt2} \otimes |f(x)\>}
\end{align*}
\onslide<5->{for some $R \subset \{0,1\}^n$}
\end{column}
\end{columns}

\end{frame}

%%%%%%%%%%%%%%%%%%%%%%%%%%%%%%%%%%%%%%%%%%%%%%%%%%%%%%%%%%%%%%%%%%%%%%%%%%%%%%%%

\begin{frame}{Quantum algorithm for Simon's problem: Analysis II}

Recall $H^{\otimes n}|x\> = \sum_{y \in \{0,1\}^n} (-1)^{x \cdot y} |y\>$

\onslide<2->{
\begin{align*}
  H^{\otimes n} \left( \frac{|x\>+|x \oplus s\>}{\sqrt2} \right)
  &\onslide<3->{= \frac{1}{\sqrt{2^{n+1}}} \sum_{y \in \{0,1\}^n} [(-1)^{x \cdot y} + (-1)^{(x \oplus s)\cdot y}]|y\>} \\
  &\onslide<4->{= \frac{1}{\sqrt{2^{n+1}}} \sum_{y \in \{0,1\}^n} (-1)^{x \cdot y} [1+(-1)^{s \cdot y}] |y\>}
\end{align*}
}

\onslide<5->{
Two cases:
\begin{itemize}
  \item if $s \cdot y = 0 \bmod 2$, $1+(-1)^{s \cdot y} = 2$
  \item<6-> if $s \cdot y = 1 \bmod 2$, $1+(-1)^{s \cdot y} = 0$
\end{itemize}
}

\bigskip

\onslide<7->{
Measuring gives a random $y$ orthogonal to $s$ (i.e., $s \cdot y = 0$)
}

\end{frame}

%%%%%%%%%%%%%%%%%%%%%%%%%%%%%%%%%%%%%%%%%%%%%%%%%%%%%%%%%%%%%%%%%%%%%%%%%%%%%%%%

\begin{frame}{Quantum algorithm for Simon's problem: Post-processing}
  
Measuring gives a random $y$ orthogonal to $s$ ($s \cdot y = 0$)

\pause
\bigskip

Repeat $k$ times, giving vectors $y_1,\ldots,y_k \in \{0,1\}^n$; solve a system of $k$ linear equations for $s \in \{0,1\}^n$:
\[
y_1 \cdot s = 0,\quad y_2 \cdot s = 0, \quad \ldots, \quad y_k \cdot s = 0
\]

\bigskip

How big should $k$ be to give a unique (nonzero) solution? \pause
\begin{itemize}
  \item Clearly $k \ge n-1$ is necessary
  \item It can be shown that $k = O(n)$ suffices
\end{itemize}

\pause
\bigskip

$O(n)$ quantum queries, $O(n^3)$ quantum gates

\pause
\bigskip

Compare to $\Omega(2^{n/2})$ classical queries (even for bounded error)

\end{frame}

%%%%%%%%%%%%%%%%%%%%%%%%%%%%%%%%%%%%%%%%%%%%%%%%%%%%%%%%%%%%%%%%%%%%%%%%%%%%%%%%

\begin{frame}{Recap}
  
We have seen several examples of quantum algorithms that outperform classical computation:
\begin{itemize}
  \item Deutsch's problem: 1 quantum query vs.\ 2 classical queries
  \item Deutsch-Jozsa problem: 1 quantum query vs.\ $2^{\Omega(n)}$ classical queries (deterministic)
  \item Simon's problem: $O(n)$ quantum queries vs.\ $2^{\Omega(n)}$ classical queries (randomized)
\end{itemize}

\bigskip

Quantum algorithms for more interesting problems build on the tools used in these examples.

\end{frame}

%%%%%%%%%%%%%%%%%%%%%%%%%%%%%%%%%%%%%%%%%%%%%%%%%%%%%%%%%%%%%%%%%%%%%%%%%%%%%%%%

\begin{frame}{Exercise: One-out-of-four search}
  
  \small

Let $f:\{0,1\}^2 \to \{0,1\}$ be a black-box function taking the value $1$ on exactly one input.  The goal is to find the unique $(x_1,x_2) \in \{0,1\}^2$ such that $f(x_1,x_2)=1$.
\begin{itemize}
  \item Write the truth tables of the four possible functions $f$.
  \item How many classical queries are needed to solve the problem?
  \item Suppose $f$ is given as a quantum black box $U_f$ acting as
  \[
    |x_1,x_2,y\> \mapsto |x_1,x_2,y \oplus f(x_1,x_2)\>.
  \]
  Determine the output of the following quantum circuit for each of the possible black-box functions $f$:
  \[
  \Qcircuit @C=1em @R=.5em {
  \lstick{|0\>} & \gate{H} & \multigate{2}{f} & \qw \\
  \lstick{|0\>} & \gate{H} & \ghost{f} & \qw \\
  \lstick{|1\>} & \gate{H} & \ghost{f} & \qw 
  }\]
  \item Show that the four possible outputs obtained in the previous part are pairwise orthogonal.  What can you conclude about the quantum query complexity of one-out-of-four search?
\end{itemize}

\end{frame}